\documentclass[11pt]{article}
\usepackage[a4paper,margin=1in]{geometry}
\usepackage[utf8]{inputenc}
\usepackage[T1]{fontenc}
\usepackage{amsmath,amssymb}
\usepackage{graphicx}
\usepackage{booktabs}
\usepackage{xcolor}
\usepackage{hyperref}

% A number not yet returned by a finished run renders loudly, so that a placeholder can never be
% mistaken for a measurement.
\newcommand{\pending}[1]{\textcolor{red}{\textbf{[PENDING: #1]}}}

\title{\textbf{Teaching to the Test: How Much of an Architecture\\ Search's Improvement Is Not Real}}
\author{Vasili Gavrilov\thanks{Independent researcher. Reference implementation:
\url{https://github.com/Anode1/SMBPANN}.}}
\date{August 2026}

\begin{document}
\fbox{\parbox{0.94\textwidth}{\textbf{RETIRED 2026-08-12. Do not extend this draft.} Its central
mechanism, that search compute converts into quality only if the selection signal is refreshed, was
falsified three ways: the same fixed replicate with four times the mutation rate beats the rotation arm;
both deterministic yardsticks plateau while both time-varying ones keep improving, including one
carrying no information; and the fixed arm's own believed score is flat, so it never learned its sample,
it stopped searching. The search underneath also loses to matched-budget random search. See
\texttt{RETIRED.md}; the measurement work continued in \texttt{\textasciitilde/bpnn}.}}
\bigskip

\maketitle

\begin{abstract}
Give a class the same practice exam every week and the marks will climb while the learning stops. An
evolutionary architecture search does something very similar, for much the same reason, and this paper
measures what it costs.

The search we were running would not do what it was asked. An earlier paper in this line had shown that
the compact filter of a convolution appears on its own once mutation acts on a shared feature rather than
on a single connection; the other half of a convolution, one filter reused at every position, it obtained
by construction and said so. We built a task where that reuse pays, gave the search a genome
free to refuse it, and waited. The search refused. A larger set of operators did not help, and neither
did twelve times the generation budget. We assumed for some time that the fault was in the search, and it
was not. The fault was in the number the search was being judged by.

A search scores each design on data held back from training, and because training is random a design can
score well by luck. Keeping the best score keeps the luck. The remedy costs one line of bookkeeping: hold
back two samples instead of one and choose using only the first. If the samples are interchangeable, then
a design chosen without consulting either must score the same on both, so any gap between the sample used
for choosing and the sample kept in reserve was manufactured by the choosing. That gap has an honest value
of zero, needs no baseline, and is denominated in the same percentage points as the result being reported.

On NAS-Bench-101, a public table of 423,624 networks each trained three separate times, an ordinary-sized
search overstates its accuracy by about half a percentage point, which is more than the improvements such
searches are usually credited with. The consequence for compute is blunter than the consequence for
honesty. Judged against one fixed held-out sample, our search reaches 94.18 after fifty generations and
94.19 after sixteen hundred, so a thirty-twofold increase in effort buys two thousandths of a point.
Alternate between two samples, one per generation, and the same compute buys two tenths. The search was
never short of compute. It was spending it to learn one sample more thoroughly.

Two repairs correct two different faults. Averaging two samples lowers the luck in every score and helps
equally at any length of search. Alternating removes the fixed target instead, so its benefit compounds,
and past roughly two hundred generations it overtakes averaging. On the small task, where the best design
is fixed in advance and verified by enumeration, repairing the yardstick lifts exact recovery of that
design from 16 runs in 240 to 42. The last finding is the one we like least and believe most: once the
yardstick was fixed, the search began beating the design we had planted for it to find.
\end{abstract}

\section{Where this continues from}

The predecessor to this work \cite{paper1} treated an evolutionary search as an instrument rather than
an optimiser. Starting from a fully connected network and charging a price for every connection, it
asked which parts of a convolution the price alone would recover. Sparse task-relevant connectivity
appeared readily, and weight sharing was adopted whenever the task rewarded recognising the same
feature in more than one place. The compact local filter turned out to depend on a detail of the
mutation operator, and that dependence was the paper's contribution: mutation acting on one connection
at a time cannot find the filter, while mutation acting on a shared offset can, because only the second
form exposes the fact that sharing lets connection count and stored-weight count move independently.

What that paper did not obtain by search was the tiling, meaning the reuse of a single filter at every
position. Its genome encoded connectivity as an offset mask, which makes every candidate translation
invariant by construction, so the reuse was imposed rather than discovered. The paper said as much.
Closing that gap is the obvious next step, and the obvious way to attempt it is to give the search a
genome in which reuse is optional and a task in which reuse pays, then see whether the search elects it.

We built that task and ran that search, and for a long time the answer looked like a clean negative. The
search settled on scattered supports rather than the compact aligned one, and neither a larger set of
mutation operators nor twelve times the generation budget changed the outcome. A negative of that shape
invites a structural explanation, and we spent a considerable amount of effort on structural
explanations. The eventual cause was elsewhere, and it was not a property of the search at all. The
quantity the search was selecting on did not measure what we were reporting, and once that was repaired
the behaviour changed. Because the fault lies in the measurement rather than in the algorithm, it is not
specific to convolutions, to evolutionary search, or to our task, which is why the bulk of this paper is
about a public benchmark rather than about our own construction.

\section{The problem, and the one change that measures it}

Architecture search in its plainest form is trial and error with a scoreboard. A design is generated,
trained, and scored on data that the training never saw. Good scores survive, are modified slightly, and
are scored again. Nothing about this is controversial, and it works.

The difficulty enters through the randomness of training. Starting weights are drawn at random and the
order of examples is shuffled, so the same design trained twice yields two different scores. Any single
score is therefore the design's real quality plus a contribution from luck, and selecting the largest
score selects for both. The luck does not average away afterwards, because the design was chosen partly
because of it. The reported number is consequently too high, and how much too high is not something a
practitioner can read off the procedure.

An everyday version of the same arithmetic shows what is going on. Give a hundred students one
examination and crown the highest scorer. That student is able, and also probably had a good
day. Set a second examination and the score falls, not because the student has deteriorated, but because
part of what the first examination selected was the good day, and good days do not repeat on demand. A
student chosen at random rather than by score shows no such fall. The difference between the chosen and
the unchosen is the entire quantity this paper measures.

\subsection{Two held-out samples instead of one}

Hold back two samples rather than one, choose using only the first, and leave the second untouched until
the end.

Suppose the two samples are interchangeable, by which we mean they are the same kind of data separated
at random, so that neither has any claim to be the more truthful. For a design chosen without consulting
either of them, the expected score on the first equals the expected score on the second. That is not an
assumption about the data or the model; it follows from the way the split was made. Now take the design
the search actually returned, which was chosen using the first sample, and compare its two scores. Any
gap between them was produced by the act of choosing, because interchangeability has already excluded
every other source.

We call that gap the inflation. It is worth stating three of its properties, because together they are
what make it usable rather than merely correct. Its honest value is exactly zero, so it requires no
control group and no baseline to be interpreted. It is expressed in the same percentage points of
accuracy as the headline result, so it can be compared directly against whatever improvement is being
claimed. And it costs one extra held-out sample, which in many settings is already available and simply
being averaged away.

\subsection{What is new, and what is not}

Almost none of the statistics here is new, and it is worth naming the prior art before claiming anything
at all. Selecting on a noisy measurement and then reporting that same measurement is a known bias, and
the special case of taking a maximum over noisy candidates is the winner's curse. That architecture
searches overfit their validation signal is likewise known within the field, where the usual advice is
to enlarge the validation set, resample it, or search less aggressively.
\pending{cite the specific NAS validation-overfitting literature, which has not yet been read and so is
not yet claimed.}

The additions are four. First, a measurement whose null value is zero by construction, which removes the
usual difficulty of deciding what an overfitting number should be compared against. Second, the size of
that number on real architecture data across four orders of magnitude of search effort, including the
observation that it exceeds the accuracy differences these papers routinely treat as decisive. Third,
evidence that two natural repairs address two different faults, one of which compounds with search
length while the other does not, so that the better repair depends on how long the search runs and the
two exchange places at a measurable point. Fourth, a controlled check on a task whose best design is
fixed in advance and verified by enumeration, where repairing the yardstick roughly triples how often
the search returns that design.

\section{How large is it? Evidence from NAS-Bench-101}\label{sec:nb101}

NAS-Bench-101 \cite{nb101} is a published table of 423,624 network designs that have already been
trained and scored, so searching over it costs a lookup instead of a training run. The property that
matters here is incidental to its original purpose: every design was trained three separate times.
Those three scores differ only through the randomness of training, which is exactly the condition for
interchangeability, so the benchmark supplies the two held-out samples at no cost.

One practical detail nearly cost us the experiment, and it illustrates how easily this structure is
discarded. The copy of the table in our own repository averages the three runs into a single number,
which is the correct table for ranking designs and the wrong one for studying the yardstick, since
averaging destroys precisely the structure the measurement depends on. We re-extracted from the original
archive, keeping every run. As a check on the extraction, the mean of our three recovered scores
reproduces the averaged table to within $10^{-4}$ percentage points on all 423,624 rows.

It is worth knowing how noisy one training run is before interpreting anything else. Across the
benchmark, retraining the same design gives scores whose spread has a median of 0.33 percentage points,
which sounds mild enough to ignore. The tail is not mild. At the 99th percentile the spread is 42.7
points, because roughly one design in a hundred sometimes trains successfully and sometimes collapses to
the accuracy of guessing. A search that trusts a single training run will occasionally be shown the
successful twin of a design that usually fails.

\subsection{The measurement}

Draw $K$ designs at random, let the search keep whichever scores best on the first held-out sample, then
look up that design's score on the sample it was never shown. Here $K$ stands for search effort.
Alongside the inflation we record two quantities that matter more to somebody actually running a search:
the true quality of the chosen design, taken as the mean of all three of its runs, and the regret,
meaning how far short of the best design that was actually present in those $K$ the search finished.

\begin{center}
\begin{tabular}{rccccc}
\toprule
 & \multicolumn{2}{c}{one held-out sample} & \multicolumn{2}{c}{two, averaged} & \\
\cmidrule(lr){2-3}\cmidrule(lr){4-5}
$K$ & inflation & regret & inflation & regret & inflation avoided \\
\midrule
$64$        & $0.291$ & $0.091$ & $0.178$ & $0.030$ & $0.113$ \\
$1{,}024$   & $0.507$ & $0.163$ & $0.342$ & $0.060$ & $0.165$ \\
$32{,}768$  & $0.636$ & $0.212$ & $0.418$ & $0.074$ & $0.218$ \\
$262{,}144$ & $0.545$ & $0.205$ & $0.382$ & $0.083$ & $0.162$ \\
\bottomrule
\end{tabular}
\end{center}

Every figure is in percentage points of accuracy over 100,000 repetitions per row, and the standard
error on the final column is 0.0015 or below, so none of these differences is in question.

The inflation is large in the terms the field itself uses. At $K = 1024$, a search size nobody would
describe as excessive, half a percentage point of the reported accuracy is not real, and papers in this
area regularly contest differences smaller than that. The inflation also grows with effort, rising
steadily from $K = 4$ to about $K = 32{,}768$, so the harder the search looks the greater the fraction
of what it finds that is luck. Past that point it eases off slightly, for a reason we prefer to state
rather than leave as an anomaly: once the sample approaches the size of the whole table the true best
designs are nearly always included, so the winner is increasingly a real winner. Averaging two
samples helps on both counts at once, removing about a third of the inflation and cutting the regret by
a factor between two and three. The second of those is the more important, because it means the repair
changes which design is taken home and not merely what is said about it.

\section{The same practice exam every week}\label{sec:repairs}

Random search has no history, so it can only show what averaging buys. A generational search has
history, and history is where the interesting fault lives.

Consider what a fixed held-out sample becomes over fifty generations. A design that happens to suit that
particular sample keeps its unearned advantage, is copied into the next generation, and is copied again,
and its descendants inherit whatever quirk of that sample favoured it. The sample stops being a test and
becomes a syllabus. This is the practice exam given every week: the marks improve, and what improves them
is familiarity with the paper rather than command of the subject. Averaging two samples does not address
this at all, since it produces a better test and then sets the same better test forever. What addresses it
is changing the paper, which costs nothing, because a second sample is already being held.

So the two repairs correct different faults, and the difference is mechanical rather than a matter of
degree. Averaging reduces the luck present in any single score, two runs having $\sqrt 2$ less spread than
one, and that benefit arrives immediately and does not grow. Alternating leaves the luck in each score
untouched, every generation still judging on one noisy run, and instead denies the population a fixed
target to converge on. Accumulation takes time, so alternation should be worth nothing in a short search
and progressively more as the search lengthens. We registered that prediction with its direction before
running the sweep.

It holds, and the magnitude at long search lengths was larger than we expected. Below is a search with a
population of 24 over the NAS-Bench-101 cell space, 20,000 runs per cell, every figure the true quality of
the design returned, in percentage points.

\begin{center}
\begin{tabular}{rcccc}
\toprule
generations & one fixed sample & alternating & averaging & alternating $-$ fixed \\
\midrule
$10$    & $94.139$ & $94.132$ & $94.238$ & $-0.006$ \\
$20$    & $94.173$ & $94.177$ & $94.290$ & $+0.004$ \\
$50$    & $94.183$ & $94.230$ & $94.313$ & $+0.048$ \\
$100$   & $94.184$ & $94.269$ & $94.318$ & $+0.085$ \\
$200$   & $94.185$ & $94.318$ & $94.314$ & $+0.133$ \\
$400$   & $94.188$ & $94.356$ & $94.319$ & $+0.169$ \\
$800$   & $94.187$ & $94.395$ & $94.316$ & $+0.208$ \\
$1600$  & $94.185$ & $94.430$ & $94.315$ & $+0.245$ \\
\bottomrule
\end{tabular}
\end{center}

The second column is the result we consider the most useful thing in this paper, and it is not about
accounting. A search judged throughout by one held-out sample arrives at 94.183 by fifty generations and
at 94.185 by sixteen hundred. Thirty-two times the compute, two thousandths of a point. The same search
with the paper changed each generation improves across the whole range and finishes a quarter of a point
ahead. Search effort converts into better designs only while the yardstick is being refreshed; spent
against a fixed one it goes into learning that sample more exactly, which the sample rewards and reality
does not.

The final column shows the accumulation directly, worth nothing at ten generations and rising steadily
thereafter. Setting the third and fourth columns side by side shows the two repairs changing places at
around two hundred generations. That crossing is not a tuning artefact, since one fault is fixed in
magnitude and the other compounds, so which repair dominates is decided by how long the search runs
rather than by any parameter we chose.

\pending{the 3200-generation cell, running, to confirm the divergence continues.}

\section{A controlled check, where the answer is known in advance}\label{sec:toy}

The benchmark establishes the size of the problem but cannot say whether repairing the yardstick lets a
search find a particular structure, since nobody knows what NAS-Bench-101's best design ought to look
like. Answering that requires a task whose answer we fixed ourselves, which returns us to the question
the first paper left open.

The task is a short strip of numbers containing a small pattern at one of ten positions, and the
structure being sought is a convolution: one small filter, reused unchanged at every position. Two
properties make the convolution the best answer rather than merely a plausible one. Training
uses only three of the ten positions while scoring uses the remainder, so a design that cannot transfer
to a position it never saw earns nothing, which is what makes reuse worth having. And each stored weight
carries a price, so a filter wider than the pattern is penalised for the width it does not need. Reuse
itself is a gene the search may switch off, which is the property the predecessor's offset-mask genome
lacked.

We verify that the target is the best answer rather than assuming it. Every one of the 127 filter shapes
in the neighbourhood of the planted one is scored under the identical objective, and the planted shape
ranks first. Two features of that check earned their inclusion through failure. It has to include shapes
that are scattered rather than solid, because scattered shapes are what the search actually explores, and
an enumeration restricted to solid ones would certify the target against the wrong family. And a rank is
a statistic rather than a fact, since an argmax over noisy scores favours whichever candidate was
luckiest, so we report the winning margin beside the standard error of the scores that produced it and
believe the ordering only when the margin clears it. Applying that rule, the planted shape is the clear
winner only when the price per weight falls within a band. Set the price too low and a wider filter
actually wins, in which case a search that declines to return our target is behaving correctly and we
would have been wrong to record a failure.

With the price inside that band we compare a search using one fixed held-out sample against a search
using both repairs, over 240 runs matched on the same random draws, and score exact recovery, meaning
that the design returned is the planted convolution itself.

\begin{center}
\begin{tabular}{lc}
\toprule
 & exact recovery, of 240 runs \\
\midrule
one fixed held-out sample        & $16$ \;($6.7\%$) \\
both repairs                     & $\mathbf{42}$ \;($17.5\%$) \\
a random design of the same size  & $1$ \;($0.4\%$) \\
\bottomrule
\end{tabular}
\end{center}

The increase is roughly threefold, at $p = 7 \times 10^{-5}$ by an exact McNemar test on the paired
runs. The final row is the control that carries the most weight, since a design with the same number of
weights placed at random essentially never coincides with the target, which rules out the possibility
that recovery is an accident of size.

\subsection{What this result is not}

The natural reading of the table above is more generous than the evidence supports, so we set out the
limits explicitly.

It is not a claim that the search found the optimum. At the same setting the repaired search outperforms
our planted convolution overall, by 0.031 on the objective at $p = 7 \times 10^{-8}$, using fewer
weights at indistinguishable transfer accuracy. There is no contradiction between that and the
enumeration, and the reconciliation is instructive: enumeration identifies the best single design
averaged over many task draws, whereas the search chooses a design for the particular draw in front of
it, and a choice made per task should beat the best fixed choice. Of the 198 runs that do not return the
planted design, 136 score better than it on their own draw.

The figure of 17.5 per cent therefore needs something other than 100 per cent to be compared against.
The planted design is the best answer for an individual draw about 13 per cent of the time, and that
figure is itself a lower bound, since picking an argmax among 127 noisy candidates loses the true winner
to lucky rivals. Set beside it, the repairs move recovery from roughly half the rate the task warrants
to slightly above it, which is a smaller claim than saying the convolution emerges and is the claim the
measurements will carry.

The controlled part of this evidence is also small by construction: one strip, ten positions, a single
hidden layer. The mechanism concerns yardsticks rather than convolutions, which is why
Section~\ref{sec:nb101} exists and why we would not rest the general claim on this section alone.

\section{Related work, and what it is owed}

The statistical phenomenon is old. Selecting on a noisy measurement and then reporting that measurement
is a known bias, and the maximum-over-candidates case is the winner's curse. We claim novelty neither
for the phenomenon nor for the remedies, but for the null of exactly zero, the magnitude on real
architecture data, and the separation of the two faults.

Weight sharing as something learned rather than imposed has a substantial history: tying parameters
through a prior \cite{nowlan}, and driving locally connected layers toward convolutions either by
regularisation or by the statistics of the data \cite{neyshabur,ingrosso}, with the reachability of
convolutional solutions under gradient descent studied directly \cite{dascoli}. The price per stored
weight used here belongs to that family and is not new.
\pending{Nowlan and Hinton 1992, Elsayed et al.\ 2020, and d'Ascoli et al.\ 2019 must be read before
these sentences stand.}

Equivariance theory supplies the reason a convolution is the right structure to plant, since a linear map
that commutes with a group action is a generalised convolution \cite{kondor,cohen}. The theorem concerns
linear maps and exact symmetry, so it provides the target rather than a proof about our search space.

The lineage of the structure itself predates its use as an architecture. Local receptive fields repeated
across positions were a finding about cortex \cite{hubel} before the first architecture built them in
using explicit position-invariant feature reuse \cite{fukushima}, which preceded the weight-sharing
formulation \cite{lecun89} that convolutional networks inherit.
\pending{confirm Hubel and Wiesel 1962 and Fukushima 1980 have been read.}

\section{Conclusion}

A search that judges candidates on held-out data is trusting a number that contains luck, and taking the
best of many such numbers retains the luck. Holding out two interchangeable samples rather than one
makes the retained luck visible as a gap whose honest value is zero. On real architecture data that gap
is about half a percentage point at ordinary search sizes, which is larger than the improvements the
field argues over, and it grows as the search works harder.

Two repairs cost almost nothing and correct different faults. Averaging two samples reduces the luck in
each score and helps equally at any search length. Alternating between them removes the fixed target
that a long search learns to please, so it helps more the longer the search runs, and past a few hundred
generations it overtakes averaging. The sharpest consequence is about compute rather than about
statistics: against a fixed yardstick a search stops improving after about fifty generations and a
sixteenfold increase in effort adds nothing, while the same search with a refreshed yardstick keeps
improving throughout.

The first paper in this line asked which pieces of a convolution a search recovers on its own and
answered it for the filter. This one takes up the piece that was left imposed and finds that the
governing constraint sits in the instrument rather than in the search, which is a less satisfying answer
about convolutions and a more useful one about method. Before asking whether a search is good enough, it
is worth checking that the number it is being judged by measures what will eventually be reported. In our
case that check amounts to one additional column in a results table, and it retired two negative results
we had already believed.

\section*{Appendix: how each result regenerates}

Every probe is a single self-contained C99 file under \texttt{validation/paper2/} in the reference
implementation, every number in this paper comes from an archived output file in the same repository, and
environment variables select the mode. Defaults are preserved throughout, and after each change a binary
built from the previous source is confirmed to reproduce the previous result table byte for byte.

\begin{center}
\begin{tabular}{ll}
\toprule
result & invocation \\
\midrule
per-run scores recovered from the archive & \texttt{./nb101\_trials nasbench\_only108.tfrecord} \\
inflation against search effort (\S\ref{sec:nb101}) & \texttt{KMIN=2 KMAX=262144 REPS=100000 ./nb101\_signal} \\
averaging against alternating (\S\ref{sec:repairs}) & \texttt{POP=24 GENS=\$G RUNS=20000 ./nb101\_search} \\
the planted-target enumeration (\S\ref{sec:toy}) & \texttt{SUBSCAN=1 SUBW=7 SEEDS=60 DRAWS=12} \\
exact recovery (\S\ref{sec:toy}) & \texttt{RAW=1 ROTPOS=\{0,1\} FEVAL=\{1,12\} LAMBDA=6} \\
paired statistics & \texttt{GROUP=2 METRICS=... ./pairstat} \\
\bottomrule
\end{tabular}
\end{center}

Paired continuous outcomes are compared with the Wilcoxon signed-rank test as primary, using an exact
null where ties permit and a tie-corrected normal approximation otherwise, always reporting which was
used, together with Hodges-Lehmann estimates and distribution-free intervals, a paired $t$ as secondary,
Holm correction across the declared family, and a minimum detectable effect so that a null result is
bounded rather than merely asserted. Paired binary outcomes use an exact McNemar test, computed as a
two-sided binomial on the discordant pairs. The tool implementing all of this is checked against
hand-computable cases before use.

Predictions for each experiment were written to a timestamped file before the corresponding run,
together with what would refute them, and scored afterwards. Two were wrong and are recorded as such: on
the small task we predicted that neither repair would suffice alone, and averaging alone did more than
alternating alone; and on the benchmark we predicted that alternating would prove inert at every search
length, which held only for short searches.

\begin{thebibliography}{9}\setlength{\itemsep}{0pt}\setlength{\parskip}{0pt}
\bibitem{paper1} V.~Gavrilov. The Imposed and Emergent Pieces of Convolution Under an Energy Budget.
 Zenodo 10.5281/zenodo.21423177, 2026.
\bibitem{nb101} C.~Ying, A.~Klein, E.~Christiansen, E.~Real, K.~Murphy, F.~Hutter. NAS-Bench-101:
 Towards Reproducible Neural Architecture Search. \emph{ICML} 2019.
\bibitem{nowlan} S.~J. Nowlan, G.~E. Hinton. Simplifying Neural Networks by Soft Weight Sharing.
 \emph{Neural Computation} 4(4), 1992.
\bibitem{neyshabur} B.~Neyshabur. Towards Learning Convolutions from Scratch. \emph{NeurIPS} 2020.
\bibitem{ingrosso} A.~Ingrosso, S.~Goldt. Data-driven emergence of convolutional structure in neural
 networks. \emph{PNAS} 119(40), 2022.
\bibitem{dascoli} S.~d'Ascoli, L.~Sagun, J.~Bruna, G.~Biroli. Finding the Needle in the Haystack with
 Convolutions. \emph{NeurIPS} 2019.
\bibitem{kondor} R.~Kondor, S.~Trivedi. On the Generalization of Equivariance and Convolution in Neural
 Networks to the Action of Compact Groups. \emph{ICML} 2018.
\bibitem{cohen} T.~Cohen, M.~Welling. Group Equivariant Convolutional Networks. \emph{ICML} 2016.
\bibitem{hubel} D.~H. Hubel, T.~N. Wiesel. Receptive fields, binocular interaction and functional
 architecture in the cat's visual cortex. \emph{J.\ Physiology} 160(1), 1962.
\bibitem{fukushima} K.~Fukushima. Neocognitron: A self-organizing neural network model for a mechanism
 of pattern recognition unaffected by shift in position. \emph{Biological Cybernetics} 36, 1980.
\bibitem{lecun89} Y.~LeCun et al. Backpropagation Applied to Handwritten Zip Code Recognition.
 \emph{Neural Computation} 1(4), 1989.
\end{thebibliography}

\end{document}
